% ============================================================================
% SECTION 1: INTRODUCTION
% ============================================================================
\section{Introduction}
\label{sec:introduction}

Deep neural networks have achieved remarkable success across various computer vision tasks, including image classification, object detection, and semantic segmentation. However, these models are vulnerable to adversarial examples---carefully crafted perturbations that are imperceptible to humans but can cause misclassification with high confidence~\cite{szegedy2014intriguing, goodfellow2015explaining}. This vulnerability poses significant security concerns for safety-critical applications such as autonomous driving, medical diagnosis, and security systems.

Two dominant architectural paradigms have emerged in computer vision: Convolutional Neural Networks (CNNs) and Vision Transformers (ViTs). CNNs leverage local connectivity and weight sharing through convolutional operations, capturing hierarchical spatial features effectively~\cite{he2016deep}. In contrast, ViTs apply the self-attention mechanism from natural language processing to image patches, enabling global context modeling from early layers~\cite{dosovitskiy2021image}. While both architectures achieve competitive performance on clean data, their behavior under adversarial attacks exhibits fundamental differences that are not yet fully understood.

Previous research has extensively studied adversarial robustness for CNNs, leading to effective defense mechanisms such as adversarial training~\cite{madry2018towards} and TRADES~\cite{zhang2019theoretically}. However, the comparative analysis of CNN and ViT robustness remains limited, with existing studies often focusing on \textit{whether} one architecture is more robust than another, rather than \textit{why} they exhibit different behaviors. Understanding the mechanistic factors underlying these differences is crucial for developing more secure deep learning systems.

% ----------------------------------------------------------------------------
% Research Questions
% ----------------------------------------------------------------------------
In this work, we move beyond simple robustness comparisons to investigate the underlying mechanisms that cause CNNs and ViTs to respond differently to adversarial perturbations. Our investigation centers on several interconnected questions: Why do adversarial perturbations transfer asymmetrically between architectures, with CNN-generated examples transferring to ViTs at significantly higher rates than the reverse? Can gradient characteristics explain the observed robustness differences? How do adversarial perturbations propagate through Vision Transformer layers, and what feature-level dynamics lead to misclassification? What design principles emerge from these mechanistic insights for building robust hybrid systems that leverage the complementary strengths of both architectures?

% ----------------------------------------------------------------------------
% Contributions
% ----------------------------------------------------------------------------
This paper makes several contributions to understanding architectural differences in adversarial vulnerability. We provide a comprehensive robustness evaluation using AutoAttack~\cite{croce2020reliable}, the gold standard benchmark, comparing adversarially trained ResNet-18 and ViT-Tiny under $\epsilon=8/255$ $L_\infty$ perturbations. Our transfer attack analysis reveals and explains the significant asymmetry in cross-architecture transferability, demonstrating that this asymmetry stems from fundamental differences in how the architectures compute gradients. Specifically, we identify that CNN gradients exhibit higher sparsity while ViT gradients show greater cross-sample alignment---properties that directly determine perturbation transferability. Furthermore, we present the first systematic layer-wise analysis of how adversarial perturbations affect ViT feature representations, revealing progressive semantic degradation that intensifies with network depth. To support reproducible research, we release our complete experimental pipeline, trained models, and analysis tools.\footnote{Code and models available at: \url{https://github.com/[anonymous]/cnn-vit-adversarial-analysis}}

% ----------------------------------------------------------------------------
% Paper Organization
% ----------------------------------------------------------------------------
We begin by reviewing related work on adversarial attacks, defenses, and architecture-specific robustness studies in Section~\ref{sec:related_work}. Section~\ref{sec:methodology} describes our experimental methodology, and Section~\ref{sec:experiments} presents results across four complementary analyses: robustness comparison, transfer attacks, gradient characteristics, and feature degradation. We discuss the implications of our findings in Section~\ref{sec:discussion} and conclude in Section~\ref{sec:conclusion}.
